\scanpage{59}
% archive foliation: 58
\begin{center}-10-\end{center}
\begin{flushright}n°385\end{flushright}
dans un voisinage ouvert de $a$. En particulier on a $\varepsilon_a=\Delta_a^0$.

6. \uline{Distributions et vecteurs tangents}.

La proposition 2 montre comment un système de coordonnées au voisinage d'un point $a$ de $X$ définit une base de l'espace des distributions de support $a$ ; cette base dépend du système de coordonnées choisi. Nous allons étudier comment varient les éléments d'ordre 1 de la base en fonction du système de coordonnées. \uline{Nous supposerons} $r\neq\unsure{\omega|0}$.

Soient $a$ un point de $X$ et $L$ l'espace vectoriel $\bar D_a^1(X)$. Comme on a $\langle\mu,f\rangle=0$ pour toute distribution $\mu\in L$ et toute fonction constante $f$, la prop.2 montre que l'application :
\[
(22)\qquad \mu\longrightarrow(\langle\mu,x_1\rangle,\ldots,\langle\mu,x_n\rangle)
\]
est un isomorphisme de $L$ sur $K^n$ quel que soit le système de coordonnées $x_1,\ldots,x_n$ sur $X$ au voisinage de $a$. On peut par contre définir une base $(D_i)_a$ de $\bar D_a^1(X)$ par $\langle(D_i)_a,x_j\rangle=\delta_{i,j}$. On a donc $(D_i)_a=\Delta_a^\alpha$ avec $\alpha=e_i$. On a de plus le résultat suivant :

LEMME 1.-- \uline{Soit} $f$ \uline{une fonction de classe} $C^r$ \uline{définie dans un voisinage ouvert} $V$ \uline{de} $a$ \uline{contenu dans} $U$, \uline{et soit} $\mu$ \uline{dans} $L$. \uline{Si} $f$ \uline{est de la forme} :
\[
(23)\qquad f(p)=\Phi(x_1(p),\ldots,x_n(p))
\qquad p\in V
\]
\uline{et si} $D_i\Phi$ \uline{est la dérivée de la fonction} $\Phi$ \uline{par rapport à la i-ème variable, on a} :
\[
(24)\qquad
\langle\mu,f\rangle=\sum_{1\leq i\leq n}\langle\mu,x_i\rangle\,D_i\Phi(x_1(a),\ldots,x_n(a))
\]
\uline{et en particulier}
\[
(25)\qquad
\langle(D_i)_a,f\rangle=D_i\Phi(x_1(a),\ldots,x_n(a))
\qquad 1\leq i\leq n.
\]

La formule (24) est évidemment satisfaite si $f=1$, ou si $f=x_i$. Supposons maintenant que le germe de $f$ en $a$ soit dans $I_a^2(X)$ ; il existe alors des fonctions
